\chapter{Bai Juyi and the Plain Style}

\section{The Mid Tang}

After the An Lushan rebellion the Tang survived, but weakened: provincial military
governors became semi-independent, the court was dominated by eunuchs and factions, and
taxes fell heavily on farmers. Mid-Tang poets responded in two opposite directions:
\begin{itemize}
  \item \textbf{Difficulty and strangeness}: Han Yu \zht{韓愈}, Meng Jiao, Li He
        \zht{李賀} (``the ghost of poetry''), who sought striking, harsh, unusual language.
  \item \textbf{Plainness and social purpose}: Bai Juyi and his friend Yuan Zhen \zht{元稹},
        who wanted poetry to be clear and to criticise misgovernment.
\end{itemize}

\section{Bai Juyi \zht{白居易} (772--846)}

Courtesy name Letian \zht{樂天} (``content with Heaven''), later styled ``the Layman of
Fragrant Mountain'' (\zht{香山居士}).

\begin{itemize}
  \item Passed the \emph{jinshi} in 800, one of the youngest in his group; held a series of
        posts as censor and remonstrator.
  \item 815: demoted to Jiangzhou (Jiujiang) for outspokenly demanding the arrest of the
        assassins of a chief minister. There he wrote ``Song of the Pipa'' (\zht{琵琶行}).
  \item Later served as governor of Hangzhou (where he repaired the West Lake dikes; the Bai
        Causeway is named in his memory) and
        Suzhou, and retired to Luoyang, writing poems about wine, music, his cat, his
        illnesses and his ageing body.
  \item Extraordinarily prolific: about 2,800 poems survive, more than any other Tang poet,
        partly because he edited and deposited copies of his own collected works in
        monasteries.
\end{itemize}

\begin{important}[Bai Juyi's poetics]
In a famous ``Letter to Yuan Zhen'' (\zht{與元九書}, 815), Bai argued that poetry should
return to the moral function of the \emph{Shijing}:
\zht{文章合為時而著，歌詩合為事而作} --- ``prose should be written for its times, songs
and poems should be composed for real events.'' He criticised the merely pretty nature
poetry of the Six Dynasties and praised the social poems of Du Fu. With Yuan Zhen he led
the \textbf{New Yuefu Movement} (\zht{新樂府運動}): poems on contemporary abuses under new,
self-invented \emph{yuefu} titles.
\end{important}

\begin{aside}
A Song anecdote says that Bai Juyi read each poem to an old peasant woman and revised
until she understood it. Whether or not it is true, it captures his reputation for
\textbf{plainness} (\zht{淺切} / \zht{平易}). Later critics sometimes found him too
plain (Su Shi: ``Yuan is frivolous, Bai is vulgar''), but he was immensely popular in his
own lifetime, including in Japan, where his collected works were a basic text of Heian
court culture (Murasaki Shikibu and Sei Shōnagon quote him).
\end{aside}

\section{A farewell poem written at sixteen}

``Grass on the Ancient Plain: Seeing Off a Friend'' \zht{賦得古原草送別} was supposedly
written for a practice examination topic (\zht{賦得} ``composed on the set topic'').

\begin{center}
\begin{tabular}{p{0.22\textwidth}p{0.27\textwidth}p{0.41\textwidth}}
\toprule
Original & Pinyin & My translation \\
\midrule
\zht{離離原上草，} & lí lí yuán shàng cǎo & Lush, lush, the grass upon the plain: \\
\zht{一歲一枯榮。} & yí suì yì kū róng & each year it withers, each year it thrives. \\
\zht{野火燒不盡，} & yě huǒ shāo bú jìn & Wildfire burns it but cannot finish it; \\
\zht{春風吹又生。} & chūn fēng chuī yòu shēng & when the spring wind blows, it grows again. \\
\zht{遠芳侵古道，} & yuǎn fāng qīn gǔ dào & Its far fragrance invades the ancient road; \\
\zht{晴翠接荒城。} & qíng cuì jiē huāng chéng & its sunlit green touches the ruined town. \\
\zht{又送王孫去，} & yòu sòng wáng sūn qù & Once more I see a young lord off; \\
\zht{萋萋滿別情。} & qī qī mǎn bié qíng & thick, thick, it is full of parting feeling. \\
\bottomrule
\end{tabular}
\end{center}

\textbf{Analysis.} Couplet 2 became proverbial for resilience (and later, revolutionary
endurance). The ending recalls the \emph{Chuci} ``Summoning the Recluse'' (``The young
lord wanders and does not return; spring grass grows thick, thick''), the same allusion
Wang Wei reversed: grass, spreading everywhere, becomes the image of feeling that follows
the traveller endlessly. A legend says that when the young Bai first visited Chang'an, the
poet Gu Kuang punned on his name (\zht{居易} ``living is easy'') --- ``rice is expensive in
Chang'an; living here is not easy!'' --- until he read this poem and said that with such
lines, living anywhere would be easy.

\section{``A Question for Liu the Nineteenth'' \zht{問劉十九}}

\begin{center}
\begin{tabular}{p{0.22\textwidth}p{0.27\textwidth}p{0.41\textwidth}}
\toprule
\zht{綠蟻新醅酒，} & lǜ yǐ xīn pēi jiǔ & Green ants on new-brewed wine, \\
\zht{紅泥小火爐。} & hóng ní xiǎo huǒ lú & a little stove of red clay. \\
\zht{晚來天欲雪，} & wǎn lái tiān yù xuě & Evening comes, the sky looks like snow: \\
\zht{能飲一杯無？} & néng yǐn yì bēi wú & could you manage a cup? \\
\bottomrule
\end{tabular}
\end{center}

\textbf{Analysis.} A twenty-character invitation. ``Green ants'' are the bubbles of
unfiltered new wine. Colour contrast (green / red / white snow implied), warmth against
cold, and a final colloquial question with the sentence-final \zht{無} (= ``or not?''),
which is nearly spoken language. The poem's charm is that it is a real note to a friend.

\section{Satire: ``The Old Charcoal Seller'' \zht{賣炭翁}}

One of the fifty \emph{New Yuefu}. Its subtitle states its purpose: ``suffering from the
palace markets'' (eunuch agents who ``bought'' goods for the palace at nominal prices).

\begin{center}
\begin{tabular}{p{0.36\textwidth}p{0.54\textwidth}}
\toprule
\zht{賣炭翁，伐薪燒炭南山中。} & The old charcoal seller cuts wood and burns charcoal in the Southern Hills. \\
\zht{滿面塵灰煙火色，} & His face is covered with dust and ash, the colour of smoke; \\
\zht{兩鬢蒼蒼十指黑。} & his temples are grey, his ten fingers black. \\
\ldots & \ldots \\
\zht{可憐身上衣正單，} & Pitiful --- the clothes on his back are thin, \\
\zht{心憂炭賤願天寒。} & yet he worries charcoal is cheap, and wishes for cold weather. \\
\bottomrule
\end{tabular}
\end{center}

He drives his cart through the snow to the market, where two palace envoys seize the
whole load ``in the name of an imperial edict'' and pay with a length of red gauze and some
silk. The poem ends there, without comment.

\begin{important}[How the satire works]
\begin{itemize}
  \item Narrative in plain language; mostly seven-character lines.
  \item The killer couplet \zht{可憐身上衣正單，心憂炭賤願天寒}: his poverty is so
        extreme that he wishes for the cold that hurts him. A single paradox shows the
        whole economic trap.
  \item The ending withholds moral judgement: the reader supplies it. (Contrast Du Fu's
        ``Officer at Shihao'', which ends similarly with a bare fact.)
\end{itemize}
\end{important}

\section{``Song of Lasting Regret'' \zht{長恨歌} (806)}

A 120-line seven-character narrative poem on the love of Emperor Xuanzong
(\zht{唐玄宗}, r.\ 712--756) and his consort Yang Guifei \zht{楊貴妃}. Bai Juyi
writes, transparently, of a ``Han emperor'' (\zht{漢皇}), a common Tang way of writing
about one's own dynasty.

\begin{center}
\begin{tikzpicture}[st/.style={draw, thick, rounded corners, align=center, text width=2.4cm,
  minimum height=1.2cm, font=\scriptsize}]
  \node[st, fill=pink!25] (a) at (0,0) {1. The emperor finds Yang; love and neglect of government};
  \node[st, fill=red!20] (b) at (3,0) {2. Rebellion; flight to Sichuan; Yang killed at Mawei};
  \node[st, fill=gray!20] (c) at (6,0) {3. Return to the palace; the emperor's grief};
  \node[st, fill=blue!15] (d) at (9,0) {4. A Daoist adept finds her in the immortals' isle};
  \draw[-{Stealth}, thick] (a) -- (b); \draw[-{Stealth}, thick] (b) -- (c);
  \draw[-{Stealth}, thick] (c) -- (d);
  \node[font=\scriptsize, align=center] at (4.5,-1.35) {history $\longrightarrow$ romance $\longrightarrow$ fantasy};
\end{tikzpicture}
\end{center}

The opening line:
\begin{center}
\zht{漢皇重色思傾國，御宇多年求不得。}\\
``The Han emperor prized beauty, longing for a woman to topple kingdoms; / for years on the
throne he sought her and did not find her.''
\end{center}

\zht{傾國} ``to topple a state'' is a Han idiom for a great beauty, and ominous: she
\emph{will}. The ending, as the immortal Yang recalls their vow at the Palace of Long Life
on the seventh night of the seventh month:

\begin{center}
\begin{tabular}{p{0.36\textwidth}p{0.54\textwidth}}
\toprule
\zht{在天願作比翼鳥，} & ``In heaven, may we be birds that fly wing to wing; \\
\zht{在地願為連理枝。} & on earth, may we be branches grown together.'' \\
\zht{天長地久有時盡，} & Heaven is long and earth is lasting, yet they will end; \\
\zht{此恨綿綿無絕期。} & this regret goes on and on, with no ending. \\
\bottomrule
\end{tabular}
\end{center}

\begin{important}[Interpreting ``Lasting Regret'']
Is it a \textbf{satire} of an emperor whose infatuation nearly destroyed the dynasty (the
opening line sets that up), or a \textbf{celebration} of romantic love that outlasts death
(the second half overwhelms the satire)? Most readers feel both. The poem was hugely
popular (Bai said singing-girls raised their price if they could recite it), inspired
plays such as the Qing drama \emph{The Palace of Lasting Life} (\zht{長生殿}), and shaped
the opening chapter of \emph{The Tale of Genji}.
\end{important}

\section*{Summary}
\begin{itemize}
  \item Mid Tang: difficulty (Han Yu, Li He) versus plainness (Bai Juyi, Yuan Zhen).
  \item Bai Juyi: ``write for real events''; New Yuefu Movement; huge, self-edited
        corpus; popular in Japan.
  \item ``Grass on the Ancient Plain'': resilience; \zht{野火燒不盡，春風吹又生}.
  \item ``Question for Liu the Nineteenth'': plainness and intimacy.
  \item ``Old Charcoal Seller'': satire by paradox (wishing for the cold).
  \item ``Song of Lasting Regret'': Xuanzong and Yang Guifei; satire or romance; vow of the
        birds and branches.
\end{itemize}
